% lecture_2.tex

\section*{Cursul II}

\begin{exercitiu}{1}
    Să se dezvolte algoritmul \textbf{GSb}. 
\end{exercitiu}

\begin{demonstratie}
    Pentru metoda \textbf{GSb} avem:
    $$
    P_{GSb} = D + U.
    $$

    Această matrice este \textbf{inversabilă} dacă și numai dacă:
    $$
    \det(D + U) \neq 0
    \quad \Longleftrightarrow \quad
    \prod_{i=1}^{n} a_{ii} \neq 0
    \quad \Longleftrightarrow \quad
    a_{ii} \neq 0, \ \forall i = 1, \dots, n.
    $$

    \vspace{1em}
    Ecuația $i$ pentru metoda \textbf{GSb} este:
    $$
    a_{i1}x_1\p{k-1} + a_{i2}x_2\p{k-1} + \dots + a_{ii}x_i\p{k} + \dots + a_{in}x_n\p{k} = b_i.
    $$

    De aici rezultă:
    \begin{equation}
    x_i\p{k} = \frac{1}{a_{ii}} 
    \left[
    b_i - 
    \sum_{j=1}^{i-1} a_{ij} x_j\p{k-1} -
    \sum_{j=i+1}^{n} a_{ij} x_j\p{k}
    \right].
    \label{eq:gsb}
    \end{equation}

    Ca algoritm, se parcurge sistemul de la \textbf{ultima} ecuație ($i = n$) până la \textbf{prima} ($i = 1$) și se calculează fiecare $x_i$ cu formula \eqref{eq:gsb}.
\end{demonstratie}

\vspace{2em}
\begin{teorema}{1}
    \text{ }\\[0.5em]
    Dacă $A \in \M_n(\R)$ este \textbf{strict diagonal dominantă}, atunci metoda \textbf{GSb} converge.
\end{teorema}

\begin{demonstratie} 
    \begin{observatie}
    Deoarece $M_{GSb}(A) = M_{GS}(A^{\mathsf T})$, rezultă că metoda \textbf{Gauss-Seidel descendentă} aplicată lui $A$ este echivalentă cu metoda \textbf{Gauss-Seidel ascendentă} aplicată lui $A^{\mathsf T}$.
    $$
    \Rightarrow\ \text{Se poate folosi demonstrația din curs pentru metoda GSa aplicată lui } A^{\mathsf T}.
    $$
    \end{observatie}
\end{demonstratie}

\begin{teorema}{2}
    \text{ }\\[0.5em]
    Dacă $A \in \M_n(\R)$ este \textbf{SPD}, atunci metoda \textbf{GSb} converge.
\end{teorema}

\begin{demonstratie} 
    \begin{observatie} În mod similar cu exercițiul anterior, $M_{GSb}(A) = M_{GS}(A^{\mathsf T})$, rezultă că metoda \textbf{Gauss-Seidel descendentă} aplicată lui $A$ este echivalentă cu metoda \textbf{Gauss-Seidel ascendentă} aplicată lui $A^{\mathsf T}$.
    $$
    \Rightarrow\ \text{Se poate folosi demonstrația din curs pentru metoda GSa aplicată lui } A^{\mathsf T}.
    $$
    \end{observatie}
\end{demonstratie}

\begin{teorema}{1 Kahan}
    \text{ }\\[0.5em]
    Pentru orice $\omega \in \R$, se are $\rho(M_{SORb}(\omega)) \geq |\omega - 1|$.
\end{teorema}

\begin{demonstratie} 
    Observăm că $M_{\mathrm{SORb}}(A,\omega)=M_{\mathrm{SOR}}(A^{\mathsf T},\omega)$ și $\rho(M)=\rho(M^{\mathsf T})$. Aplicând teorema lui Kahan pentru $M_{\mathrm{SOR}}(A^{\mathsf T},\omega)$, obţinem afirmaţia.
\end{demonstratie}

\begin{corolar}
    Metoda \textbf{SORb} nu converge pentru $\omega \in \R \setminus (0,2)$.
\end{corolar}

\begin{demonstratie} 
    Din teoremă, $\rho(M_{\mathrm{SORb}})\ge |1-\omega|\ge 1$ când $\omega\notin(0,2)$.
\end{demonstratie}

\begin{teorema}{2 Ostrowski-Reich}
    \text{ }\\[0.5em]
    Dacă $A \in \M_n(\R)$ este \textbf{SPD}, atunci metoda \textbf{SORb} converge $\Leftrightarrow \omega \in (0,2)$.
\end{teorema}

\begin{demonstratie} 
    \begin{observatie}
        Fiindcă $P_{\mathrm{SOR}}(A^{\mathsf T},\omega)=P_{\mathrm{SORb}}(A,\omega)$ și 
        $N_{\mathrm{SOR}}(A^{\mathsf T},\omega)=N_{\mathrm{SORb}}(A,\omega)$, rezultă
        $M_{\mathrm{SOR}}(A^{\mathsf T},\omega)=M_{\mathrm{SORb}}(A,\omega)$.
    \end{observatie}
    Prin urmare, demonstrația din curs pentru \textbf{SOR} aplicată lui $A^{\mathsf T}$ e identică la \textbf{SORb} aplicată lui $A$. 
\end{demonstratie}

\begin{teorema}{1}
    \text{ }\\[0.5em]
    (i) Dacă $A$ este simetrică $\Rightarrow$ matricile $P_{SSOR}(\omega)$ și $N_{SSOR}(\omega)$ sunt simetrice.\\[0.5em]
    (ii) Dacă $A$ este simetrică și $D$ pozitiv definită $\Rightarrow$ $M_{SSOR}(\omega)$ are valori proprii reale nenegative.\\[0.5em]
    (iii) Dacă $A$ este simetrică și $D$ pozitiv definită $\Rightarrow$ $\big[P_{SSOR}(\omega) + P_{SSOR}(\omega)^T - A\big]$ este SPD $\Leftrightarrow \omega \in (0,2)$.
\end{teorema}

\begin{demonstratie} 
    $A=D+L+U$, $D^{\mathsf T}=D$, $U=L^{\mathsf T}$ (deoarece $A$ este simetrică). Rezultă că:
    \begin{gather*}
        P_{SSOR}(\omega)=\frac{1}{\omega(2-\omega)}(D+\omega L)\,D^{-1}\,(D+\omega L)^{\mathsf T},
        N_{SSOR}(\omega)=\frac{1}{\omega(2-\omega)}\big((1-\omega)D-\omega L\big)\,D^{-1}\,\big((1-\omega)D-\omega L\big)^{\mathsf T}.
    \end{gather*}
    \noindent \textbf{(i)} \\
    Fiecare matrice este de forma $X\,D^{-1}X^{\mathsf T}$ $\Rightarrow$ matricile $P_{SSOR}(\omega)$ și $N_{SSOR}(\omega)$ sunt simetrice.

    \noindent \textbf{(ii)} 
    $$
    M_{SSOR}(\omega) = (D + \omega U)^{-1} D (D + \omega L)^{-1} \big[ (1 - \omega) D - \omega L \big] D^{-1} \big[ (1 - \omega) D - \omega U \big].
    $$
    $A = D + L + U$ este simetrică $\Rightarrow$ $U = L^{\mathsf T}$ și $D = D^{\mathsf T}$.  
    Notăm:
    $$
    X := (D + \omega L)^{-1} \big[ (1 - \omega) D - \omega L \big] D^{-1/2}.
    $$
    Atunci:
    $$
    M_{SSOR}(\omega) = D^{-1/2} X^{\mathsf T} X D^{1/2}.
    $$

    Matricea $X^{T}X$ este simetrică și semipozitiv definită, deci valorile proprii ale ei sunt reale și nenegative.

    Deoarece $M_{\text{SSOR}}$ este similară cu $X^{T}X$ (deoarece $D^{1/2}$ este inversabilă, iar $D^{1/2}M_{\text{SSOR}}D^{-1/2}=X^{T}X$) $\Rightarrow$ $M_{\text{SSOR}}$ are aceleași valori proprii ca $X^{T}X$, deci reale și nenegative.

    \noindent \textbf{(iii)} 

    \begin{align*}
        P_{SSOR}(\omega) &=\frac{1}{\omega(2-\omega)}(D+\omega L)\,D^{-1}(D+\omega U),\\[0.4em]
        P_{SSOR}(\omega)^{\mathsf T} &=\frac{1}{\omega(2-\omega)}(D+\omega U)\,D^{-1}(D+\omega L),\\[0.6em]
        P_{SSOR}(\omega)+P_{SSOR}(\omega)^{\mathsf T} &=\frac{1}{\omega(2-\omega)}\Big[(D+\omega L)D^{-1}(D+\omega U) +(D+\omega U)D^{-1}(D+\omega L)\Big],\\[0.6em]
        &=\frac{1}{\omega(2-\omega)} \Big[(D+\omega L)(I+\omega D^{-1}U) +(D+\omega U)(I+\omega D^{-1}L)\Big],\\[0.6em]
        &=\frac{1}{\omega(2-\omega)} \Big[\,2D+2\omega(L+U)+\omega^2\!\big(LD^{-1}U+UD^{-1}L\big)\Big].
    \end{align*}

    \begin{align*}
        P_{SSOR}(\omega)+P_{SSOR}(\omega)^{\mathsf T}-A
        &=\frac{1}{\omega(2-\omega)}
        \Big[\,2D+2\omega(L+U)+\omega^2(LD^{-1}U+UD^{-1}L)\\[-0.2em]
        &\hspace{8.7em}-\omega(2-\omega)D-\omega(2-\omega)(L+U)\Big]
    \end{align*}

    \begin{align*}
        &=\frac{1}{\omega(2-\omega)}
        \Big[(2-\omega(2-\omega))D
        +(2\omega-\omega(2-\omega))(L+U)
        +\omega^2\!\big(LD^{-1}U+UD^{-1}L\big)\Big].
    \end{align*}

    Folosind $U=L^{\mathsf T}$ și identitatea
    $$
    LD^{-1}U+UD^{-1}L=(L+U)-D,
    $$
    obținem:
    \begin{align*}
        P_{SSOR}(\omega)+P_{SSOR}(\omega)^{\mathsf T}-A &= \frac{1}{\omega(2-\omega)} \Big[(2-2\omega+\omega^2)D+\omega^2\big((L+U)-D\big)+\omega^2(L+U)\Big]\\[0.3em]
        &=\frac{1}{\omega(2-\omega)}\Big[(2-2\omega)D\Big] = \frac{2-\omega}{\omega}\,D.
    \end{align*}

    $$
    \boxed{P_{SSOR}(\omega) + P_{SSOR}(\omega)^{\mathsf T} - A = \dfrac{2 - \omega}{\omega}\,D.}
    $$

    Cum $D$ este PD, rezultă că termenul din stânga este SPD ddacă $\dfrac{2-\omega}{\omega}>0$, adică \,$\omega\in(0,2)$.
\end{demonstratie}

\begin{teorema}{2 Ostrowski-Reich}
    \text{ }\\[0.5em]
    Dacă $A \in \M_n(\R)$ este \textbf{SPD}, atunci metoda \textbf{SSOR} converge $\Leftrightarrow \omega \in (0,2)$.
\end{teorema}

\begin{demonstratie} 
    Dacă SSOR converge, atunci $\rho(M_{SSOR})<1$.
    Prin criteriul \emph{Householder-John}:
    $$
    P_{SSOR}(\omega)+P_{SSOR}(\omega)^{\mathsf T}-A\ \text{este SPD}.
    $$
    Din demonstrația anterioară,
    $$
    P_{SSOR}+P_{SSOR}^{\mathsf T}-A=\frac{2-\omega}{\omega}\,D,
    $$
    iar $D$ este PD (deoarece $A$ este SPD). Așadar $(2-\omega)/\omega>0$, adică $\omega\in(0,2)$.

    Dacă $\omega\in(0,2)$, atunci $\dfrac{2-\omega}{\omega}\,D$ este SPD, deci $P_{SSOR}+P_{SSOR}^{\mathsf T}-A$ este SPD. Aplicând criteriul \emph{Householder-John} $\Rightarrow$ $\rho(M_{SSOR})<1$ $\Rightarrow$ metoda SSOR converge.

\end{demonstratie}